\section{Results}
\label{sec:results}

% Results report numbers. No mechanism, no conservation implication, no
% "suggesting that". Every one of those belongs in the Discussion.

\subsection{Connectivity change, \studyperiod}
\TODO{PLAIN-LANGUAGE WRITING GUIDE.
What happened: I compared all 253 possible pairs of cores. Overall, the modeled cost
of moving between cores changed only a little from 2012 to 2024. The middle change was
-1.224 percent. A negative value means the model became slightly easier to move
through; it does not mean that cheetahs actually moved more. Of the 253 pairs, 92
became harder and 161 became easier. The largest increase was +13.1 percent for the
link between cores 33 and 38.
What the map adds: outside the cores, the places with the highest modeled current were
mostly in the same locations through time. Their overlap scores were 0.91 to 0.93,
where 1 would mean perfect overlap.
What not to claim: do not call this an overall connectivity loss, observed movement,
or a country-by-country result. The analysis did not assign each core pair to a
country. For H1, report that 18 of the 45 selected links became harder and 27 became
easier, but that change was not meaningfully related to baseline betweenness
($\rho \approx 0.01$). Therefore H1 was mixed rather than supported. Use
\figref{fig:change}. When you point to a figure in your own sentences, type
\texttt{\textbackslash figref\{fig:change\}} rather than ``Figure'' by hand: it prints
``Figure'' in the capstone and ``Fig.'' in the journal layouts, which is how Springer writes it.\\
H1, TWO MORE NUMBERS YOU CAN USE (checked 5 October 2026; same answer, stated the way the protocol worded H1). The $\rho \approx 0.01$ above mixes the links that got harder with the ones that got easier. Asking the protocol's question directly: (1) Did the 18 harder links have higher betweenness than the 27 easier ones? Slightly, but not meaningfully (Mann--Whitney one-sided $p = 0.36$). (2) Among the 18 harder links, did the bigger losses fall on the more important links? No; if anything the reverse ($\rho = -0.22$, not significant). One thing that DID match: the losses were fairly concentrated, with the 5 hardest-hit of the 18 links carrying 46 percent of the total increase. So H1's ``minority of links'' half holds, and its ``high-betweenness links'' half does not.\\
FIGURE PLACEMENT: the study-area figure just below shows the extent and cores, which is setup, not a result. Move it to Methods, near Study area.}

\begin{figure}[htbp]\centering
  \includegraphics[width=\textwidth]{Figure_0_study_area_and_core_groups.pdf}
  \caption{Study extent, fixed historical range cores, and core groups used to define
  the modeled connectivity network. Core groups are an analysis construct and do not
  imply separate populations.}
  \label{fig:studyarea}
\end{figure}
\FloatBarrier

\subsection{Priority links and sensitivity checks}
\TODO{PLAIN-LANGUAGE WRITING GUIDE.
What happened: 32 of the 45 selected links passed the tested sensitivity screens
and remained in the main priority group. Network importance was ranked separately
using betweenness. The other 13 links were kept as
uncertain links where more field information would be useful. The weighting test kept
34 links, and the fence test removed two more. The vegetation test did not remove any
additional links from the remaining group.
What it means: the main group is not a small handful of links; it contains most of the
network. That only partly matches the original expectation that priorities would be
limited to a small subset.
What not to claim: these links are modeled connections between cores, not confirmed
animal travel routes. Use Table~\ref{tab:top20}.}

\begin{figure}[htbp]\centering
  \includegraphics[width=\textwidth]{Figure_5_modelled_links_and_transportation.pdf}
  \caption{Modeled core-to-core links with transportation and available fence context.
  Lines represent network links used to define least-cost analyses; they are not
  observed animal paths.}
  \label{fig:networkcontext}
\end{figure}
\FloatBarrier

% Sensitivity evidence explains the priority-link screening above.
\TODO{PLAIN-LANGUAGE WRITING GUIDE FOR THE SENSITIVITY TESTS.
Main lesson: the total modeled cost of a connection was usually fairly stable, but
the exact line taken by the least-cost path was less stable. Of 180 route-and-year
comparisons, 49 had less than 80 percent overlap within 1 km. This affected 21 of the
32 priority links. The route between cores 17 and 24 had no overlap in 2020.
Vegetation assumptions: 43 of 45 links kept the same direction of change under all
three vegetation weightings. However, the size of the change grew as vegetation was
given more weight: the middle absolute change increased from about 3.6 to 9.3 percent.
Network-design assumption: changing how many neighbors each core could connect to
changed the number of links and changed which links ranked highest. A bridge is a link
whose removal splits the network into separate pieces; link 17--24 was the only bridge
when three nearest neighbors were used.
What not to claim: there is no completed cell-by-cell agreement map, and no single
ranking should be presented as certain. State that H3 was only partly evaluated: a
 stable group existed, but it contained a majority of the selected links and the full
planned sensitivity grid was not completed.}

\TODO{PLAIN-LANGUAGE WRITING GUIDE FOR THE INDEPENDENT TRACKING CHECK (UPDATED 5 OCTOBER 2026 WITH THE STRICTER FOLLOW-UPS).\\
Give this its own short subsection; it answers a different question from the priority links. Numbers only here; meaning goes in Discussion. Lead with the step check (Table~\ref{tab:limpopo_steps}), because it is the stricter test. Put the first comparison (Table~\ref{tab:limpopo}) after it as background, or move that table to \suppB.\\[0.3em]
1. THE MAIN RESULT (Table~\ref{tab:limpopo_steps}, panel A). Across 549 day-to-day moves, the place a cheetah actually went had lower modeled resistance than the moves it could have made instead. The real move beat 56 percent of its alternatives (95 percent interval 51 to 61 percent), where 50 would mean no difference, and seven of the nine animals leaned that way. Averaged along the straight line between locations the result was almost the same (55 percent).\\
2. CURRENT DID NOT HOLD UP. The real move beat 53 percent of alternatives on modeled current, but the interval (47 to 59) includes 50, and only five of nine animals leaned that way. Say this plainly. It also means the 65 percent for current in the first comparison mostly reflected where the cats lived, not the choices they made each day.\\
3. ALL NINE SURFACES (panel B). Every surface gave 55 to 56 percent, so the result does not depend on which weighting you picked. The four other fence treatments give exactly the same values as the main surface because no mapped fence lies within 20~km of any track.\\
4. WHICH LAYERS CARRY IT (panel C). Road density did most of the work: on its own it reached 55 percent, about as much as the full surface, and leaving it out dropped the full surface to 54 percent with an interval that just touches 50. Livestock, slope, and night lights each added a small lean (51 to 52 percent). Built-up area did nothing here because there is almost none in this district. Vegetation on its own leaned slightly the wrong way (48 percent, interval 43 to 52, which includes 50). Report that honestly; do not hide it.\\
5. ROAD CROSSINGS (panel D). 26 percent of real moves crossed a mapped road, against 35 percent of the alternatives; eight of nine animals crossed less than expected. Major roads were crossed too rarely by anyone (1 against 2 percent) to say much.\\
6. THE FIRST COMPARISON (Table~\ref{tab:limpopo}), as background: against random points across the whole area, cheetah locations had lower resistance 60 percent of the time (7 of 9 animals) and higher current 65 percent of the time (6 of 9). Widening the area by 20~km and using the earlier model gave the same picture.\\[0.3em]
HOW TO DESCRIBE THE STRENGTH: consistent in direction, modest in size, and mostly carried by roads. 56 percent is a real lean, not a strong one. Say ``consistent with'' the model, not ``confirmed.'' Report the intervals and animal counts, not p-values.\\[0.3em]
WHAT NOT TO CLAIM: that the model is validated; that it works across the whole region (one district in South Africa); that it says anything about 2020 or 2024 (the tracks end in 2008); or that fences were tested (they were not).}

\begin{table}[htbp]\centering
  \caption{\TODO{Final caption, in your own words: day-to-day moves of nine free-roaming cheetahs (Limpopo, 2003--2008) compared with moves of the same length the animals could have made instead, on the 2012 surfaces. Panels: main surface, all nine surfaces, single input layers, road crossings.}}
  \label{tab:limpopo_steps}
  \resizebox{\textwidth}{!}{\begin{tabular}{llrcr}
\toprule
 & Comparison & Share of alternatives & 95\% interval & Animals \\
 & & the real move beat & (resampling animals) & above 0.50 \\
\midrule
\multicolumn{5}{l}{\textit{A. Main surface, 549 day-to-day moves, 20 alternatives each}}\\
 & Resistance at the end point (lower expected) & 0.56 & 0.51--0.61 & 7 of 9\\
 & Resistance along the straight-line path & 0.55 & 0.51--0.59 & 7 of 9\\
 & Current at the end point (higher expected) & 0.53 & 0.47--0.59 & 5 of 9\\
 & Current along the straight-line path & 0.52 & 0.47--0.58 & 5 of 9\\
\midrule
\multicolumn{5}{l}{\textit{B. All nine surfaces, resistance at the end point}}\\
 & Main surface & 0.56 & 0.51--0.61 & 7 of 9\\
 & Equal human-pressure weights & 0.56 & 0.51--0.61 & 7 of 9\\
 & Infrastructure emphasis & 0.56 & 0.50--0.61 & 6 of 9\\
 & Vegetation weight 10\% & 0.56 & 0.50--0.61 & 7 of 9\\
 & Vegetation weight 40\% & 0.55 & 0.51--0.60 & 6 of 9\\
 & Four other fence treatments$^{a}$ & 0.56 & 0.51--0.61 & 7 of 9\\
\midrule
\multicolumn{5}{l}{\textit{C. Single input layers, resistance at the end point}}\\
 & Road density & 0.55 & 0.50--0.62 & 6 of 9\\
 & Livestock & 0.52 & 0.49--0.56 & 5 of 9\\
 & Terrain (slope) & 0.52 & 0.50--0.54 & 4 of 9\\
 & Night lights & 0.51 & 0.51--0.52 & 8 of 9\\
 & Built-up area$^{b}$ & 0.50 & 0.50--0.50 & 3 of 9\\
 & Vegetation & 0.48 & 0.43--0.52 & 3 of 9\\
 & Main surface without road density & 0.54 & 0.50--0.57 & 7 of 9\\
\midrule
\multicolumn{5}{l}{\textit{D. Road crossings along the straight-line path (fewer expected)}}\\
 & Any mapped road (26\% of real moves crossed, 35\% of alternatives) & 0.54 & 0.52--0.57 & 8 of 9\\
 & Major roads only (1\% against 2\%) & 0.51 & 0.50--0.51 & 5 of 9\\
\bottomrule
\end{tabular}}
\\[0.5em]
\noindent\TODO{Final note, in your own words. Each real day-to-day move (at least 1~km) is compared with 20 moves the same cheetah could have made instead: same start, a step length drawn from that animal's own moves, random direction. The score is the share of those alternatives the real move beat; 0.50 means no difference. Each animal counts once, and the 95\% interval comes from resampling the nine animals. $^{a}$No mapped fence lies within 20~km of any track, so all five fence treatments give identical values here; these data cannot test the fence part of the model. $^{b}$Built-up area barely varies in this district, so it cannot separate real and alternative moves.}

\end{table}

\begin{table}[htbp]\centering
  \caption{\TODO{Final caption, in your own words: modeled resistance and current at
  daily locations of nine free-roaming cheetahs (Limpopo, 2003--2008) compared with
  random available points from the same landscape, on the 2012 surfaces. Call it the broader, less strict comparison.}}
  \label{tab:limpopo}
  \resizebox{\textwidth}{!}{\begin{tabular}{llrrrrr}
\toprule
Surface & Measure & Cheetah fixes & Available & Expected direction & 20-km buffer & Animals\\
 & & (median) & (median) & (share) & (share) & (of 9)\\
\midrule
Final model (vegetation-balanced, documented fences) & Resistance (lower expected) & 2.23 & 2.45 & 0.60 & 0.60 & 7\\
 & Current (higher expected) & 0.156 & 0.145 & 0.65 & 0.62 & 6\\
Earlier model (pressure and terrain only, no fences) & Resistance (lower expected) & 2.03 & 2.34 & 0.61 & 0.61 & --\\
 & Current (higher expected) & 0.143 & 0.128 & 0.65 & 0.63 & --\\
\bottomrule
\end{tabular}}
\\[0.5em]
\noindent\TODO{Final note, in your own words: the comparison uses 1,006 daily locations from nine free-roaming cheetahs (2003--2008) against 10,060 random points from the same landscape, on the 2012 surfaces. ``Expected direction (share)'' is the chance that a random cheetah location had lower resistance (or higher current) than a random available point; 0.50 means no difference. ``20-km buffer'' repeats this with a wider available area. ``Animals'' counts how many of the nine had a median in the expected direction. No location fell inside a core. Add one sentence saying this is the broader comparison, and that in the stricter step check (Table~\ref{tab:limpopo_steps}) the resistance lean held but the current lean did not.}

\end{table}

\TODO{PLAIN-LANGUAGE WRITING GUIDE FOR THE SAME-ARCHIVE OCCURRENCE CHECK.
Keep this, but present it after the tracking check as the weaker of the two, because it is not independent.
What happened: after thinning records to one 10-km cell, the early period had 413
occupied cells and the later period had 450. Their overlap score was 0.239, showing
that the two sets were quite different. Later records occurred in slightly higher
modeled resistance, and a larger share fell inside the fixed cores. Almost no records
were close to the priority paths: zero early cells and three later cells.
Why this is a weak check: the same occurrence archive helped create the cores, so this
is not an independent test. There are also no occurrence records after 2016, so it
cannot check the 2020 or 2024 results. Report it honestly as a limited comparison, not
as validation.}

\TODO{Protocol reporting only: the planned cell-level monitoring index was not built and H4 was not tested. Explain this once in \suppA, not as a repeated priority-link result.}

\subsection{Protected-area coverage of priority connectivity}
\TODO{PLAIN-LANGUAGE WRITING GUIDE.
Why the cores were removed: the model forces current into and out of the cores, so the
cores automatically receive very high current. Including them would make protection
look better without telling us how well the land between cores is protected.
What happened outside the cores: 31.26 percent of the comparison area is protected.
Only 23.27 percent of the area above the 80th current percentile was protected in
2024. Coverage rose to 25.54 percent above the 90th percentile, 33.45 percent above the
95th, and 44.63 percent above the 99th. In simple terms, the very strongest
concentrations often reached protected areas, while the broader connecting landscape
around them was less protected.
What not to claim: do not report a significance p-value because it changed when the
size of the spatial blocks changed. Do not call modeled current observed movement. Use
Table~\ref{tab:pa}. State that H2 was not supported under its full rule because the
direction changed across current thresholds and the significance result was not
stable.}

%------------------------------------------------------------------ figure inserts
\begin{figure}[htbp]\centering
  \includegraphics[width=\textwidth]{Figure_2B_current_change_2012_2024.pdf}
  \caption{Modeled current change between 2012 and 2024.}
  \label{fig:change}
\end{figure}
\FloatBarrier

\begin{figure}[htbp]\centering
  \includegraphics[width=\textwidth]{Figure_2A_modelled_current_2024.pdf}
  \caption{Modeled current concentration for 2024, with the fixed core network shown
  for orientation.}
  \label{fig:current2024}
\end{figure}
\FloatBarrier

\begin{table}[htbp]\centering
  \caption{Covariates, native resolution, source years, and static/dynamic treatment.}
  \label{tab:covariates}
  \small\begin{tabular}{>{\raggedright\arraybackslash}p{0.22\textwidth}>{\raggedright\arraybackslash}p{0.24\textwidth}>{\raggedright\arraybackslash}p{0.13\textwidth}>{\raggedright\arraybackslash}p{0.34\textwidth}}
\toprule
Covariate & Source & Treatment & Temporal role / notes\\
\midrule
Built-up fraction & Global Human Settlement Layer (GHSL) & Dynamic & 2010, 2015, 2020, 2025 products matched to snapshots\\
Nighttime radiance & Visible Infrared Imaging Radiometer Suite (VIIRS), VNL V2 & Dynamic & 2013, 2016, 2020, 2024 products\\
Tree cover & Moderate Resolution Imaging Spectroradiometer (MODIS), MOD44B & Dynamic & Annual snapshot fields\\
Non-tree vegetation & Moderate Resolution Imaging Spectroradiometer (MODIS), MOD44B & Dynamic & Annual snapshot fields\\
Bare cover & Moderate Resolution Imaging Spectroradiometer (MODIS), MOD44B & Dynamic & Annual snapshot fields\\
Road density & Global Roads Inventory Project, version 4 (GRIP4) & Static & Spatial pressure/context\\
Livestock density & Gridded Livestock of the World, version 4 (GLW4) & Static & Combined cattle, goats, and sheep exposure\\
Terrain / slope & Shuttle Radar Topography Mission (SRTM) & Static & Slope-derived resistance\\
Hydrography & Hydrological data and maps based on Shuttle Elevation Derivatives at multiple Scales (HydroSHEDS) & Static & Water/context mask\\
Veterinary fences & World Wide Fund for Nature (WWF), Kavango--Zambezi Transfrontier Conservation Area (KAZA); Kruger references & Scenario & Sensitivity/context, not assumed impassable\\
\bottomrule
\end{tabular}
\\[0.5em]
\parbox{\textwidth}{\footnotesize\emph{Note:} Product shorthand retained in the source names is defined at first use: VNL = VIIRS nighttime lights; V2 = product version 2; MOD44B = MODIS vegetation-continuous-fields product.}
\\[0.5em]
\noindent\TODO{THREE GAPS IN THIS TABLE, in plain terms.\\
1. The reader cannot tell which layers actually went into the resistance surface. Your final model used built-up, roads, livestock, and nighttime lights (human pressure), vegetation, and slope. Hydrography is listed as ``Water/context mask,'' but it was not one of the weighted inputs. Add a column or a note saying ``in resistance: yes/no'' so nobody assumes water was a barrier.\\
2. Categorical land cover (MCD12Q1) is not listed, but Methods says it was used for masking and description. Add it as ``Context only, never differenced.''\\
3. Protected areas (WDPA, August 2026) are not listed, but they drive a whole Results section. Add them as ``Overlay for reporting only, not a resistance input.''}

\end{table}

\begin{table}[htbp]\centering
  \caption{Ten highest-ranked modeled priority links shown from the 32-link primary-priority set.}
  \label{tab:top20}
  \resizebox{\textwidth}{!}{\begin{tabular}{rrrrrrr}
\toprule
From & To & Rank & Betweenness & Type & Distance (km) & Change (\%)\\
\midrule
38 & 40 & 1 & 0.409 & KNN+MST & 120.6 & +4.78\\
24 & 38 & 2 & 0.316 & KNN & 202.3 & +7.44\\
17 & 24 & 3 & 0.300 & KNN+MST & 208.0 & -5.78\\
30 & 40 & 4 & 0.238 & KNN & 169.2 & +3.55\\
27 & 41 & 5 & 0.103 & KNN+MST & 196.8 & -6.45\\
41 & 45 & 6 & 0.090 & KNN & 183.6 & -5.50\\
33 & 38 & 7 & 0.079 & KNN & 108.1 & +13.12\\
9 & 17 & 8 & 0.079 & KNN+MST & 109.0 & -5.26\\
26 & 27 & 9 & 0.079 & KNN+MST & 221.8 & -5.18\\
40 & 43 & 10 & 0.077 & KNN & 159.6 & -1.56\\
\bottomrule
\end{tabular}}
\parbox{\textwidth}{\footnotesize\emph{Note:} The table shows the ten highest-ranked members of the 32-link primary-priority set; links are modeled structural network relationships, not observed animal routes.}

\end{table}

\begin{table}[htbp]\centering
  \caption{Protected-area coverage of modeled high-current classes.}
  \label{tab:pa}
  \resizebox{\textwidth}{!}{\begin{tabular}{rrrrrr}
\toprule
Year & Current percentile & Cells & Protected (\%) & Background (\%) & Difference (pp)\\
\midrule
2012 & 80 & 418,295 & 24.23 & 31.26 & -7.03\\
2012 & 90 & 209,148 & 26.31 & 31.26 & -4.95\\
2012 & 95 & 104,574 & 35.57 & 31.26 & +4.31\\
2012 & 99 & 20,915 & 44.17 & 31.26 & +12.92\\
2024 & 80 & 418,292 & 23.27 & 31.26 & -7.99\\
2024 & 90 & 209,146 & 25.54 & 31.26 & -5.72\\
2024 & 95 & 104,573 & 33.45 & 31.26 & +2.19\\
2024 & 99 & 20,915 & 44.63 & 31.26 & +13.37\\
\bottomrule
\end{tabular}}
\\[0.5em]
\noindent\TODO{Final note: values use the outside-core, positive-current domain. Percentage-point differences are relative to the matched background. Do not report a permutation $p$-value because inference was sensitive to the assumed spatial block size.}

\end{table}

\begin{table}[htbp]\centering
  \caption{Overview of completed modeled-connectivity results.}
  \label{tab:priority}
  \begin{tabular}{p{0.53\textwidth}p{0.16\textwidth}p{0.25\textwidth}}
\toprule
Metric & Value & Interpretation\\
\midrule
Final modeled-current snapshots & 4 (2012, 2016, 2020, 2024) & 253/253 focal-region pairs connected each year\\
Fixed density cores & 23 & Historical baseline, fixed across snapshots\\
Selected neighboring links & 45 & Core-pair comparison set\\
Primary-priority links & 32 & Stable under weight/fence screens\\
All-pair median effective-resistance change & -1.224\% & 2012--2024\\
Selected-link median change & -1.557\% & 2012--2024\\
Primary-priority median change & -2.618\% & 2012--2024\\
Largest all-pair increase & +13.119\% & Link 33--38\\
Largest all-pair decrease & -8.243\% & Link 43--45\\
\bottomrule
\end{tabular}

\end{table}
